\documentclass[10pt]{article}

% -------------------------
% PAQUETES
% -------------------------
\usepackage[spanish]{babel}
\usepackage[utf8]{inputenc}
\usepackage[T1]{fontenc}
\usepackage{geometry}
\usepackage{setspace}
\usepackage{graphicx}
\usepackage{amsmath}
\usepackage{booktabs}
\usepackage{caption}
\usepackage{hyperref}
\usepackage{apacite}
\usepackage{microtype}

% -------------------------
% CONFIGURACIÓN DE PÁGINA (según pauta)
% -------------------------
\geometry{
 top=1.5cm,
 left=1.5cm,
 right=1.5cm,
 bottom=2cm
}
\setstretch{1.0} % interlineado simple
\setlength{\parindent}{0pt}
\setlength{\parskip}{0.5em}

% -------------------------
% FIGURAS / TABLAS (rótulo inferior, numeración consecutiva)
% -------------------------
\captionsetup{font=small, labelfont=bf}
\renewcommand{\figurename}{Fig.}
\renewcommand{\tablename}{Tabla}

% -------------------------
% IDENTIFICACIÓN DE GRUPO (arriba, sin ocupar una hoja)
% -------------------------
\title{\textbf{Título del trabajo / tarea}}

% Mismos estudiantes que antes
\author{
Natalay Huaiquinao, Carlos Ulloa\\
\texttt{nhuaiquinao2021@alu.uct.cl, carlos.ulloa2020@alu.uct.cl}
}

% Ajuste del encabezado según lo pedido (Semestre I-2026)
\date{\vspace{-0.6em}Tarea 1, INFO1184, Semestre I-2026}

% -------------------------
\begin{document}

\maketitle
\vspace{-1.0cm}

% -------------------------
% 2) RESUMEN (150 a 200 palabras)
% -------------------------
\begin{abstract}
Este resumen debe contener entre 150 y 200 palabras. Indique el objetivo de la tarea, el contexto del problema, y el enfoque seguido para resolverlo (por ejemplo: análisis, modelamiento, diseño metodológico paso a paso y validación con evidencias). Resuma los resultados principales, el aporte de la solución (qué mejora o qué permite concluir) y mencione brevemente supuestos o limitaciones relevantes. El texto debe ser claro, conciso y coherente, de modo que el lector comprenda el contenido general del informe sin necesidad de leer el documento completo. Evite definiciones extensas y detalles secundarios; priorice qué se hizo, cómo se hizo y qué se obtuvo. Si se incluyen referencias, utilícelas con moderación.
\end{abstract}

% -------------------------
% 3) INTRODUCCIÓN
% -------------------------
\section{Introducción}
Presente el contexto, la problemática y el enunciado del trabajo (preguntas/caso/investigación/ejercicios). Declare el objetivo general y los objetivos específicos. Delimite el alcance y describa la estructura del documento: Fundamentos de teoría, Desarrollo, Discusión y reflexión, Conclusiones, Referencias y Anexo.

% -------------------------
% 4) FUNDAMENTOS DE TEORÍA
% -------------------------
\section{Fundamentos de teoría}
Explique los conceptos, definiciones y fundamentos necesarios para comprender la solución. Integre revisión bibliográfica y cite en norma APA usando \texttt{apacite}. Ejemplos: \citeA{autorEjemplo2020} o (\citeauthor{autorEjemplo2022}, \citeyear{autorEjemplo2022}). Aclare la notación, el lenguaje técnico utilizado y la forma de representación del problema (modelos, supuestos, variables).

% -------------------------
% 5) DESARROLLO DEL TRABAJO
% -------------------------
\section{Desarrollo del trabajo}
Desarrolle la solución de forma procedimental, paso a paso, mostrando evidencia y verificación (pruebas, casos, métricas, resultados). En el informe principal no incluya código fuente si existe una implementación computacional; describa el procedimiento y resultados. Si corresponde, agregue scripts o fragmentos importantes solo en el Anexo.

\subsection{Ejemplo de figura (numeración consecutiva + cita si es extraída)}
\begin{figure}[h]
    \centering
    % Reemplace por su archivo real
    \includegraphics[width=0.85\textwidth]{ejemplo.png}
    \caption{Título de la figura. (\citeauthor{autorEjemplo2020}, \citeyear{autorEjemplo2020})}
    \label{fig:ejemplo}
\end{figure}

\subsection{Ejemplo de tabla (numeración consecutiva + cita si es extraída)}
\begin{table}[h]
\centering
\begin{tabular}{lll}
\toprule
Variable & Descripción & Valor/Ejemplo \\
\midrule
V1 & Texto & 123 \\
V2 & Texto & 456 \\
\bottomrule
\end{tabular}
\caption{Título de la tabla. (\citeauthor{autorEjemplo2020}, \citeyear{autorEjemplo2020})}
\label{tab:ejemplo}
\end{table}

\subsection{Ejemplo de ecuación (numeración a la derecha)}
Si la ecuación proviene de una fuente, menciónela en el texto (p. ej., \citeA{autorEjemplo2022}).
\begin{equation}
E = mc^2
\label{eq:energia}
\end{equation}

% -------------------------
% 6) DISCUSIÓN Y REFLEXIÓN
% -------------------------
\section{Discusión y reflexión}
Interprete los resultados y explique el aporte de la solución en base a la metodología/método empleado. Justifique evidencias (qué se evaluó, cómo se verificó, por qué es consistente). Analice limitaciones, fuentes de error, riesgos, y mejoras posibles. Conecte explícitamente la discusión con lo presentado en el Desarrollo.

% -------------------------
% 7) CONCLUSIONES (por integrante)
% -------------------------
\section{Conclusiones}

\subsection*{Conclusión de Natalay Huaiquinao}
Conclusión individual: qué se logró, evidencia clave, aprendizaje y mejoras propuestas.

\subsection*{Conclusión de Carlos Ulloa}
Conclusión individual: qué se logró, evidencia clave, aprendizaje y mejoras propuestas.

% -------------------------
% 8) REFERENCIAS (APA)
% -------------------------
\section{Referencias}
\bibliographystyle{apacite}
\bibliography{referencias}

% -------------------------
% 9) ANEXO + BITÁCORA (no cuenta para 5-6 hojas)
% -------------------------
\appendix
\section{Anexo}
Use esta sección para profundizar el desarrollo, añadir procedimientos relevantes y/o incluir fragmentos importantes de script. También puede incluir un enlace a drive/repositorio (si aplica) para el código completo.

\subsection{Bitácora de trabajo}
\begin{itemize}
    \item Día 1 (AAAA-MM-DD): actividades realizadas, decisiones, división de tareas, evidencias.
    \item Día 2 (AAAA-MM-DD): avances, problemas encontrados, cómo se abordaron, pendientes.
    \item Día 3 (AAAA-MM-DD): validación, resultados finales, preparación de entrega.
\end{itemize}

\end{document}